
\section{Source Code}
\begin{lstlisting}[language=Python, caption={One-Crossing Switch Algorithm for Gauss Codes}]

import csv
from collections import defaultdict
import pandas as pd


def tokenize_gauss_code(code):
    """
    Convert a compact Gauss code string into a list of tokens.
    Example:
        'O1-O2-U1-O3+U2-U3+'
    becomes:
        ['O1-', 'O2-', 'U1-', 'O3+', 'U2-', 'U3+']
    """
    code = code.strip()
    if not code:
        return []

    tokens = []
    current = code[0]

    for char in code[1:]:
        if char in ("O", "U"):
            tokens.append(current)
            current = char
        else:
            current += char

    tokens.append(current)
    return tokens


def parse_tokens_to_dataframe(tokens, knot_name, label):
    """
    Convert Gauss code tokens into a DataFrame with columns:
    knot, label, step, kind, index, sign
    """
    rows = []

    for step, token in enumerate(tokens):
        kind = token[0]
        index = int(token[1:-1])
        sign_char = token[-1]
        sign = 1 if sign_char == "+" else -1

        rows.append({
            "knot": knot_name,
            "label": label,
            "step": step,
            "kind": kind,
            "index": index,
            "sign": sign
        })

    return pd.DataFrame(rows)


def switch_crossing(df_knot, crossing_index):
    """
    Perform a crossing switch:
    1. Swap O and U
    2. Reverse sign
    """
    df_new = df_knot.copy()
    mask = df_new["index"] == crossing_index

    df_new.loc[mask, "kind"] = df_new.loc[mask, "kind"].map({"O": "U", "U": "O"})
    df_new.loc[mask, "sign"] *= -1

    return df_new


def dataframe_to_gauss_string(df_knot):
    """
    Convert DataFrame back to Gauss code string.
    """
    df_sorted = df_knot.sort_values("step")
    tokens = []

    for _, row in df_sorted.iterrows():
        sign_char = "+" if row["sign"] > 0 else "-"
        tokens.append("{}{}{}".format(row["kind"], row["index"], sign_char))

    return " ".join(tokens)


OUTPUT_FILENAME = "all_versions_strings.csv"
label_counts = defaultdict(int)

with open(INPUT_FILENAME, "r") as fin, open(OUTPUT_FILENAME, "w", newline="") as fout:
    writer = csv.writer(fout)
    writer.writerow(["knot", "label", "version", "switched_crossing", "gauss_string"])

    for line in fin:
        entry = line.strip()

        if not entry:
            continue
        if entry.startswith("----") or entry.startswith("Number of knots"):
            continue

        parts = entry.split(maxsplit=1)
        if len(parts) < 2:
            continue

        label, code = parts

        label_counts[label] += 1
        knot_name = "{}_{}".format(label, label_counts[label])

        tokens = tokenize_gauss_code(code)
        df_knot = parse_tokens_to_dataframe(tokens, knot_name, label)

        original_string = dataframe_to_gauss_string(df_knot)
        writer.writerow([knot_name, label, "orig", "", original_string])

        for crossing_index in sorted(df_knot["index"].unique()):
            df_variant = switch_crossing(df_knot, crossing_index)
            variant_string = dataframe_to_gauss_string(df_variant)

            writer.writerow([
                knot_name,
                label,
                "switch_{}".format(crossing_index),
                crossing_index,
                variant_string
            ])

\end{lstlisting}

\begin{lstlisting}[language=Python, caption={Simplification Algorithm Using $R1$ and $R2$ Moves}]

import pandas as pd
import numpy as np
from dataclasses import dataclass


# ----------------------------
# Data structure for tokens
# ----------------------------
@dataclass(frozen=True)
class Tok:
    ou: str      # 'O' or 'U'
    cid: int     # crossing index
    sign: str    # '+' or '-'


# ----------------------------
# Parsing utilities
# ----------------------------
def parse_gauss_string(s):
    if pd.isna(s):
        return []
    parts = str(s).replace(",", " ").split()
    return [Tok(p[0], int(p[1:-1]), p[-1]) for p in parts]


def toks_to_string(toks):
    return " ".join("{}{}{}".format(t.ou, t.cid, t.sign) for t in toks)


# ----------------------------
# R1 detection and application
# ----------------------------
def find_r1_index(toks):
    n = len(toks)
    for i in range(n):
        a = toks[i]
        b = toks[(i + 1) % n]

        if a.cid == b.cid and a.sign == b.sign and a.ou != b.ou:
            return i
    return None


def apply_r1_once(toks, i):
    j = (i + 1) % len(toks)

    if i < j:
        return toks[:i] + toks[i + 2:]
    else:
        return toks[1:-1]


# ----------------------------
# R2 detection and application
# ----------------------------
def find_r2_index(toks):
    n = len(toks)
    if n < 4:
        return None

    for i in range(n):
        block = [toks[(i + k) % n] for k in range(4)]
        ids = [t.cid for t in block]

        # Must be pattern a b a b with distinct crossings
        if ids[0] != ids[2] or ids[1] != ids[3] or ids[0] == ids[1]:
            continue

        # Overs adjacent AND unders adjacent
        if not (
            block[0].ou == block[1].ou and
            block[2].ou == block[3].ou and
            block[0].ou != block[2].ou
        ):
            continue

        a0, a2 = block[0], block[2]
        b1, b3 = block[1], block[3]

        # Each crossing must have one O and one U with same sign
        if a0.ou == a2.ou or b1.ou == b3.ou:
            continue
        if a0.sign != a2.sign or b1.sign != b3.sign:
            continue

        # Crossings must have opposite signs
        if a0.sign == b1.sign:
            continue

        return i

    return None


def apply_r2_once(toks, i):
    remove = {(i + k) % len(toks) for k in range(4)}
    return [t for k, t in enumerate(toks) if k not in remove]


# ----------------------------
# Simplification procedure
# ----------------------------
def simplify_r1_r2_steps(toks):
    r1_steps = []
    r2_steps = []
    current = toks

    # Apply R1 moves
    while True:
        i = find_r1_index(current)
        if i is None:
            break
        current = apply_r1_once(current, i)
        r1_steps.append(toks_to_string(current))

    # Apply R2 moves (with interleaved R1)
    while True:
        i = find_r2_index(current)
        if i is None:
            break

        current = apply_r2_once(current, i)
        r2_steps.append(toks_to_string(current))

        # After each R2, re-apply R1
        while True:
            j = find_r1_index(current)
            if j is None:
                break
            current = apply_r1_once(current, j)
            r1_steps.append(toks_to_string(current))

    return r1_steps, r2_steps, toks_to_string(current)


# ----------------------------
# Apply to dataset
# ----------------------------
df = pd.read_csv(INPUT_FILENAME)

r1_all = []
r2_all = []
final_all = []

for s in df["gauss_string"]:
    toks = parse_gauss_string(s)
    r1, r2, final = simplify_r1_r2_steps(toks)

    r1_all.append(r1)
    r2_all.append(r2)
    final_all.append(final)


# ----------------------------
# Expand steps into columns
# ----------------------------
max_r1 = max(len(x) for x in r1_all)
max_r2 = max(len(x) for x in r2_all)

out = df.copy()

for k in range(max_r1):
    out["R1_step{}".format(k + 1)] = [
        x[k] if len(x) > k else np.nan for x in r1_all
    ]

for k in range(max_r2):
    out["R2_step{}".format(k + 1)] = [
        x[k] if len(x) > k else np.nan for x in r2_all
    ]

out["simplified_final"] = final_all

out.to_csv("knots_with_R1_R2_steps.csv", index=False)

\end{lstlisting}

\begin{lstlisting}[language=Python, caption={Two-Crossing Switch Algorithm for Gauss Codes}]

import csv
from collections import defaultdict
from itertools import combinations
import pandas as pd


def tokenize_gauss_code(code):
    """
    Convert a compact Gauss code string into a list of tokens.

    Example:
        'O1-O2-U1-O3+U2-U3+'
    becomes:
        ['O1-', 'O2-', 'U1-', 'O3+', 'U2-', 'U3+']
    """
    code = code.strip()
    if not code:
        return []

    tokens = []
    current = code[0]

    for char in code[1:]:
        if char in ("O", "U"):
            tokens.append(current)
            current = char
        else:
            current += char

    tokens.append(current)
    return tokens


def parse_tokens_to_dataframe(tokens, knot_name, label):
    """
    Convert Gauss code tokens into a DataFrame with columns:
    knot, label, step, kind, index, sign
    """
    rows = []

    for step, token in enumerate(tokens):
        kind = token[0]
        index = int(token[1:-1])
        sign_char = token[-1]
        sign = 1 if sign_char == "+" else -1

        rows.append({
            "knot": knot_name,
            "label": label,
            "step": step,
            "kind": kind,
            "index": index,
            "sign": sign
        })

    return pd.DataFrame(rows)


def switch_crossings(df_knot, crossing_indices):
    """
    Perform crossing switches on multiple crossings by:
    1. Swapping O and U
    2. Reversing the sign
    """
    df_new = df_knot.copy()
    crossing_set = set(crossing_indices)
    mask = df_new["index"].isin(crossing_set)

    df_new.loc[mask, "kind"] = df_new.loc[mask, "kind"].map({"O": "U", "U": "O"})
    df_new.loc[mask, "sign"] *= -1

    return df_new


def dataframe_to_gauss_string(df_knot):
    """
    Convert DataFrame back to a spaced Gauss code string.
    """
    df_sorted = df_knot.sort_values("step")
    tokens = []

    for _, row in df_sorted.iterrows():
        sign_char = "+" if row["sign"] > 0 else "-"
        tokens.append("{}{}{}".format(row["kind"], row["index"], sign_char))

    return " ".join(tokens)


OUTPUT_FILENAME = "all_versions_strings_switch2.csv"
label_counts = defaultdict(int)

with open(INPUT_FILENAME, "r") as fin, open(OUTPUT_FILENAME, "w", newline="") as fout:
    writer = csv.writer(fout)

    writer.writerow([
        "knot",
        "label",
        "version",
        "switched_crossing_1",
        "switched_crossing_2",
        "gauss_string"
    ])

    for line in fin:
        entry = line.strip()

        if not entry:
            continue
        if entry.startswith("----") or entry.startswith("Number of knots"):
            continue

        parts = entry.split(maxsplit=1)
        if len(parts) < 2:
            continue

        label, code = parts

        label_counts[label] += 1
        knot_name = "{}_{}".format(label, label_counts[label])

        tokens = tokenize_gauss_code(code)
        df_knot = parse_tokens_to_dataframe(tokens, knot_name, label)

        crossing_ids = sorted(df_knot["index"].unique())

        original_string = dataframe_to_gauss_string(df_knot)
        writer.writerow([knot_name, label, "orig", "", "", original_string])

        for i, j in combinations(crossing_ids, 2):
            df_variant = switch_crossings(df_knot, (i, j))
            variant_string = dataframe_to_gauss_string(df_variant)

            writer.writerow([
                knot_name,
                label,
                "switch_{}_{}".format(i, j),
                i,
                j,
                variant_string
            ])

\end{lstlisting}

